\section{Identification of the physical critical point}
\label{sec:criticality}

The matched temperature $\beta_*$ was selected by the scale map. To
identify it with the magnetization threshold, we first construct local
FK barriers at $\beta_*$ on both lattices. These barriers rule out
percolation at the matched temperature. Above $\beta_*$, the dual odds
decrease; a decision-tree argument then makes dual connections rare
enough to force primal percolation. This is the local sharpness
argument of~\cite[Section~8]{R}, applied to the exact dual family
rather than to a distributionally self-dual model.

Write $B_r(x)=\{y:|y-x|_\infty\le r\}$ on the lattice under
consideration; radii may be rounded to integers. Let
$\phi_{\beta,\omega}^{\mathrm w}$ be the primal infinite-volume
wired FK measure. It is obtained by decreasing limits of finite
wired-box laws on local increasing events, and the limit does not
depend on the exhaustion. The analogous notation with a star refers
to the dual odds at the same physical parameter $\beta$.

\subsection{Annular barriers at the matched temperature}

\begin{lemma}[Local FK barriers on both lattices]
\label{lem:fk-barriers}
There exist $a_0>1$ and $c_0>0$ with the following property, separately
for the primal and dual lattices. Let $\mathcal A(x,r)$ be the annular
graph between the lattice layers of radii $r$ and $a_0r$, and let
$\mathcal R(x,r)$ be the event of an ordinary open path between
these layers. Call its bond environment good if
\begin{equation}
 \phi_{\mathcal A(x,r),\beta_*,\omega}^{\mathrm w}
        (\mathcal R(x,r))\le1-c_0,
 \label{eq:good-annulus}
\end{equation}
using the odds of that lattice and wiring both boundary components
together. Goodness depends only on the original bonds in the annulus
and an incident layer, and
\begin{equation}
 \sup_x\Prob\bigl(\mathcal A(x,r)\text{ is not good}\bigr)
                         \longrightarrow0\qquad(r\to\infty).
 \label{eq:annulus-good-probability}
\end{equation}
For a good annulus, the same crossing upper bound holds with any
less wired boundary condition and after any decrease in its edge odds.
In particular,
\begin{equation}
       \E_\rho\phi_{\beta_*,\omega}^{\mathrm w}
                      (0\leftrightarrow\infty)=0.
 \label{eq:no-percolation-at-matching}
\end{equation}
\end{lemma}

\begin{proof}
Choose $a_0$ large enough to place two disjoint annular subbands
between the boundary layers, with positive-width free buffers around
each. In the pure critical Ising model, simultaneous plus and minus
spin circuits, one in each subband, have probability at least
$2c_0>0$, uniformly over spin pins beyond those buffers. This is
\cite[Section~2, Proposition ``Pure buffered RSW consequences'']{S}:
the two prescribed colors have disjoint buffered rectangle
constructions. Fix the subbands with enough clearance that a small
erosion of the buffers preserves this lower bound.

Proposition~\ref{prop:local-comparison}, applied to these two circuit
tests, transfers the lower bound to the quenched model, with
environmental probability tending to one. Proposition~\ref{prop:matching}
makes this comparison available for both the primal and dual arrays
at $\beta_*$. Apply it to Edwards--Sokal spins for the annular FK
measure with its fully wired boundary cluster pinned plus. Every
ordinary open FK path has constant spin. A radial path would have
to meet both separating circuits, and hence could not be open.
The triangulation used for the spin tests causes no exception:
ordinary nearest-neighbor edges cannot cross a triangulated circuit
without meeting one of its vertices. Thus the simultaneous circuit
event implies the complement of $\mathcal R(x,r)$ and proves
\eqref{eq:good-annulus} outside an event of probability tending to
zero.

Defining goodness by the original finite-volume FK probability in
\eqref{eq:good-annulus} removes the auxiliary random choices of the
scale construction. It is a measurable function of the stated local
bond set. The all-exterior-pin scope of
Proposition~\ref{prop:local-comparison} gives
\eqref{eq:annulus-good-probability} uniformly in placement.
Monotonicity in the odds and in boundary wiring gives the remaining
finite-volume assertion.

To deduce \eqref{eq:no-percolation-at-matching}, choose concentric
annuli with a fixed large geometric separation, so that their full
bond neighborhoods are disjoint. Start at a radius beyond which
the probability of a bad annulus is at most $1/2$. For annulus $i$
put $a_i(\omega)=1-c_0\1_{\{i\text{ is good}\}}$. In any containing
finite wired box $G$, condition on all edges off this annulus. Its
conditional law has a boundary partition dominated by full wiring;
therefore
\[
 \phi_{G,\beta_*,\omega}^{\mathrm w}
       (\mathcal R_i\mid\text{edges off annulus }i)
                \le a_i(\omega).
\]
Successively applying these conditional bounds shows that the
probability of crossing $\ell$ such annuli is at most
$\prod_{i=1}^{\ell}a_i(\omega)$. No independence of FK crossings
is asserted here. Independence of the original bond neighborhoods
does give
\[
 \E_\rho\prod_{i=1}^{\ell}a_i(\omega)
       =\prod_{i=1}^{\ell}\E_\rho a_i(\omega)
       \le(1-c_0/2)^\ell.
\]
The same local bounds pass to the wired infinite-volume limit.
Every infinite path crosses all the annuli, so letting
$\ell\to\infty$ proves the claim.
\end{proof}

\subsection{Dual arms above the matched temperature}

The local comparison supplies no rate in
\eqref{eq:annulus-good-probability}. We first obtain a useful rate
by using disjoint bond sets in many separated annuli at each center.
The resulting environmental exceptional probability decays faster
than every power of the box size. Outside this exceptional set a
layer exploration has small revealment, and the weighted derivative
strengthens polynomial arm decay to a stretched exponential bound.

\begin{lemma}[Rapid annealed decay of dual arms]
\label{lem:dual-arms}
For every $\beta'>\beta_*$ and $A>0$ there is a finite constant
$C_{\beta',A}$ such that, on the dual lattice,
\begin{equation}
 \E_\rho\phi_{B_{2N}(x),\beta',\omega}^{\mathrm w,*}
  \bigl(x\leftrightarrow\partial B_N(x)
                    \text{ inside }B_N(x)\bigr)
                     \le C_{\beta',A}N^{-A}.
 \label{eq:dual-arm-bound}
\end{equation}
The estimate is uniform in the center. It also holds for the same
ordinary path event under any finite or infinite-volume dual FK law
whose domain contains $B_{2N}(x)$, with arbitrary boundary conditions
outside that box.
\end{lemma}

\begin{proof}
Fix $x$. For every $z\in B_{N+1}(x)$, choose
$\ell_N\asymp\log N$ concentric annuli of the type in
Lemma~\ref{lem:fk-barriers}, with inner radii between $N^{1/4}$
and a sufficiently small multiple of $N^{1/2}$. Choose their
geometric spacing so that their full bond neighborhoods are disjoint
and all lie inside $B_{N^{1/2}}(z)$. Let $\eta_N$ be an upper bound
for their bad-environment probabilities, uniformly over centers and
the selected radii. Equation~\eqref{eq:annulus-good-probability}
gives $\eta_N\to0$.

Independence of the annuli at one fixed center shows that
\[
 \Prob\bigl(\text{at least }\ell_N/2
                       \text{ annuli at }z\text{ are bad}\bigr)
                  \le2^{\ell_N}\eta_N^{\ell_N/2}.
\]
Let $\mathcal G_N$ be the event that fewer than half the annuli
are bad at every $z\in B_{N+1}(x)$. A union bound gives
\begin{equation}
 \Prob(\mathcal G_N^c)
       \le C N^2 2^{\ell_N}\eta_N^{\ell_N/2}
       =o(N^{-A})\qquad\text{for every }A>0.
 \label{eq:amplified-environment-bound}
\end{equation}
Indeed, $\ell_N$ is bounded above and below by positive multiples
of $\log N$, while $\log(1/\eta_N)\to\infty$. There is no
independence assertion between the annulus families at different
centers.

In an environment in $\mathcal G_N$, multiply the conditional
annulus bounds exactly as in the preceding proof. Since the dual
odds decrease as $\beta$ increases, the same environment gives
\begin{equation}
 \phi_{B_{2N}(x),\beta,\omega}^{\mathrm w,*}
 \bigl(z\leftrightarrow\partial B_{N^{1/2}}(z)
                 \text{ inside }B_{N^{1/2}}(z)\bigr)
       \le(1-c_0)^{\ell_N/2}
       \le C N^{-\alpha}
 \label{eq:uniform-short-dual-arms}
\end{equation}
for a fixed $\alpha>0$, all $z\in B_{N+1}(x)$, and every
$\beta\in[\beta_*,\beta']$. For large $N$ all the small boxes
are contained in $B_{2N}(x)$; their induced boundary conditions are
covered by full annular wiring.

Fix such an environment and let
\[
 q_\omega(\beta)=\phi_{B_{2N}(x),\beta,\omega}^{\mathrm w,*}
  \bigl(x\leftrightarrow\partial B_N(x)
                    \text{ inside }B_N(x)\bigr).
\]
Choose an integer $i$ uniformly between $N/4$ and $N/2$, and explore
all ordinary open clusters in $B_N(x)$ that meet the layer
$\partial B_i(x)$. This exploration determines the event defining
$q_\omega$: any path from $x$ to $\partial B_N(x)$ meets the
chosen layer, and all clusters meeting it have been completely
explored inside $B_N(x)$. This adapts the layer exploration in
\cite[Lemma~3.2 and its proof]{DCRT}.

An edge is queried only when an endpoint lies on, or has an ordinary
open path to, the chosen layer. For each fixed endpoint $z$, the
fraction of choices of $i$ whose layer comes within distance
$2N^{1/2}$ of $z$ is at most $C N^{-1/2}$. For any other choice,
the connection to the layer forces an ordinary arm from $z$ to
$\partial B_{N^{1/2}}(z)$. Thus
\eqref{eq:uniform-short-dual-arms} gives
\begin{equation}
 \max_e\delta_e
       \le C\bigl(N^{-1/2}+N^{-\alpha}\bigr)
       \le C N^{-d},\qquad d=\min\{1/2,\alpha\}>0.
 \label{eq:layer-revealment}
\end{equation}
The tree queries only edges in $B_N(x)$, but its law is the full
wired measure on $B_{2N}(x)$. In particular, it never follows a
wired boundary identification. Edges outside $B_N(x)$ have zero
revealment.

All finite odds are positive, so $0<q_\omega(\beta)<1$.
Apply \eqref{eq:fk-osss} to this finite wired measure. In the
dual version of \eqref{eq:weighted-derivative}, the coefficients
of $-q_\omega'$ are $a_e(\beta)$, uniformly bounded below on
$[\beta_*,\beta']$. Covariances for edges outside the explored
box remain nonnegative. Consequently
\begin{equation}
 -q_\omega'(\beta)
       \ge c_{\beta'} N^d q_\omega(\beta)(1-q_\omega(\beta)),
       \qquad \beta\in[\beta_*,\beta'].
 \label{eq:dual-differential-inequality}
\end{equation}
Moreover \eqref{eq:uniform-short-dual-arms} at $z=x$ gives
$q_\omega(\beta_*)\le C N^{-\alpha}$. Integrating the derivative
of the logit, for all sufficiently large $N$,
\[
 \frac{q_\omega(\beta')}{1-q_\omega(\beta')}
 \le\frac{q_\omega(\beta_*)}{1-q_\omega(\beta_*)}
       e^{-c_{\beta'}(\beta'-\beta_*)N^d}
 \le e^{-c_{\beta'}(\beta'-\beta_*)N^d}.
\]
Averaging this estimate on $\mathcal G_N$ and using
\eqref{eq:amplified-environment-bound} on its complement proves
\eqref{eq:dual-arm-bound}. In a larger graph, condition outside
$B_{2N}(x)$ and dominate the induced boundary partition by full
wiring. This proves the final assertion, which then passes to
infinite-volume limits.
\end{proof}

\subsection{Percolation and plus magnetization}

\begin{proposition}[Physical criticality]
\label{prop:physical-criticality}
For the original iid bond law and every sufficiently small fixed
$\varepsilon>0$, the temperature in
Proposition~\ref{prop:matching} satisfies
\begin{equation}
                    \beta_* = \beta_c^\rho(\varepsilon).
 \label{eq:physical-critical-point}
\end{equation}
The environmental mean of the infinite-volume plus magnetization
is zero for $0\le\beta\le\beta_*$ and strictly positive for every
$\beta>\beta_*$.
\end{proposition}

\begin{proof}
By Lemma~\ref{lem:fk-barriers} and monotonicity,
\[
 \E_\rho\phi_{\beta,\omega}^{\mathrm w}
          (0\leftrightarrow\infty)=0,
                    \qquad0<\beta\le\beta_*.
\]
Fix $\beta'>\beta_*$. Finite planar FK duality sends the primal
odds to $2/w_e(\beta')$, exactly the dual family considered in
Lemma~\ref{lem:dual-arms}. In each buffered interior box the
induced dual boundary condition is dominated by full wiring. Thus
the lemma bounds ordinary closed-dual arms under finite primal
wired laws and, by taking local limits, under
$\phi_{\beta',\omega}^{\mathrm w}$. Only this local comparison
is needed; no distributional self-duality is used.

A closed-dual circuit surrounding $B_m(0)$ whose maximal distance
from the origin lies in $[R,2R)$ has diameter at least $cR$.
It therefore contains an ordinary dual arm reaching distance $c'R$
from some dual vertex in $B_{2R+1}(0)$. Sum
\eqref{eq:dual-arm-bound} over the $O(R^2)$ possible vertices,
using the bound with its surrounding collar, and then over dyadic
$R$ comparable with or larger than $m$. For any $A>2$ this gives
\begin{equation}
 \E_\rho\phi_{\beta',\omega}^{\mathrm w}
   \bigl(\text{a closed-dual circuit surrounds }B_m(0)\bigr)
       \le C_{\beta',A}\sum_{j\ge0}(2^j m)^{2-A}
                           \longrightarrow0.
 \label{eq:surrounding-dual-circuits}
\end{equation}

If every open cluster meeting $B_m(0)$ were finite, their finite
union would contain the box and would be connected as a vertex
set: each of its vertices is joined by an open path to the box,
and the box itself is nearest-neighbor connected. Every outgoing
edge of this union is closed. Its outer dual boundary therefore
contains a surrounding closed-dual circuit: one may first fill the
finite complementary components, after which all edges toward the
infinite complementary component are still closed. The complement of the
event in \eqref{eq:surrounding-dual-circuits} forces an infinite
open cluster to meet $B_m(0)$. For all sufficiently large $m$ this
has positive probability under the joint environment--FK law.
That law is translation invariant: the environment is iid and
the wired infinite-volume specification is covariant under
translations. Hence
\[
 0<\E_\rho\phi_{\beta',\omega}^{\mathrm w}
             (B_m(0)\leftrightarrow\infty)
 \le |B_m(0)|\,
       \E_\rho\phi_{\beta',\omega}^{\mathrm w}
             (0\leftrightarrow\infty).
\]
This proves positive expected wired percolation for every
$\beta'>\beta_*$.

It remains to identify this order parameter with the precise plus
magnetization used in the definition of $\beta_c^\rho(\varepsilon)$.
Fix an environment and $\beta>0$. In the finite spin box, include
all bonds with at least one endpoint in that box and wire the exterior
endpoints, pinning their cluster plus. Edwards--Sokal gives
\[
 \langle\sigma_0\rangle^+_{B_m,\beta,J}
       =\phi^{+,m}_{\beta,\omega}
               (0\leftrightarrow\text{wired exterior boundary}).
\]
For these square boxes the boundary probability is exactly
$u_{m+1}(\omega)$, where $u_n$ is the root-to-boundary probability
in the ordinary wired box $B_n$. Indeed, completing the exterior
endpoints to $\partial B_{m+1}$ adds only wired vertices and
wired-to-wired edges, which do not affect this probability.
The monotone limit of $u_n$ exists. To identify it with percolation
in the local infinite-volume limit, first use wired domination:
\[
 \phi_{\beta,\omega}^{\mathrm w}(0\leftrightarrow\infty)
                         \le u_m(\omega)
 \quad\text{for every }m.
\]
Conversely, for every fixed $r<m$, the root-to-boundary event in
$B_m$ implies an ordinary connection from $0$ to $\partial B_r$.
Taking $m\to\infty$ for this local event and then $r\to\infty$
shows
\[
 \lim_m u_m(\omega)
       \le\lim_r\phi_{\beta,\omega}^{\mathrm w}
                    (0\leftrightarrow\partial B_r)
       =\phi_{\beta,\omega}^{\mathrm w}(0\leftrightarrow\infty).
\]
The equality with $u_{m+1}$ gives the same limit for the exterior-plus
convention. Thus, in each environment,
\begin{equation}
 \lim_{m\to\infty}\langle\sigma_0\rangle^+_{B_m,\beta,J}
       =\phi_{\beta,\omega}^{\mathrm w}(0\leftrightarrow\infty).
 \label{eq:magnetization-percolation}
\end{equation}
All variables lie in $[0,1]$, so environmental expectations preserve
this equality. At $\beta=0$ the magnetization is zero directly.
The vanishing and positivity already proved establish
\eqref{eq:physical-critical-point} from the definition of the
critical inverse temperature.
\end{proof}
